% ============================================================
\chapter{Int\'egrale de Lebesgue}
\label{ch:integrale}
% ============================================================

Nous voici au c\oe ur de l'\'edifice~: l'\emph{int\'egrale de Lebesgue}. Tandis que Riemann d\'ecoupait l'axe des $x$ en petits intervalles et sommait les aires des rectangles, Lebesgue a eu l'id\'ee g\'eniale de d\'ecouper l'axe des $y$~: regrouper les points o\`u la fonction prend des valeurs proches, mesurer ces ensembles de niveaux, et sommer. Cette approche, pr\'esent\'ee dans sa th\`ese de 1902, produit une int\'egrale qui co\"incide avec celle de Riemann lorsque cette derni\`ere existe, mais qui s'applique \`a une classe de fonctions incomparablement plus vaste --- et surtout, qui poss\`ede des th\'eor\`emes de convergence (Fatou, convergence domin\'ee, convergence monotone) d'une puissance sans \'equivalent.

\section{Int\'egrale des fonctions simples}

\begin{definition}[Int\'egrale d'une fonction simple positive]
Soit $(X, \mathcal{A}, \mu)$ un espace mesur\'e et $\varphi = \sum_{k=1}^m c_k
\mathbf{1}_{A_k}$ une fonction simple mesurable positive ($c_k \geq 0$). On d\'efinit~:
\[
\int_X \varphi\, d\mu = \sum_{k=1}^m c_k\, \mu(A_k),
\]
avec la convention $0 \cdot (+\infty) = 0$.
\end{definition}

\begin{proposition}[Propri\'et\'es de l'int\'egrale des fonctions simples]
\label{prop:integrale_simples}
Soient $\varphi, \psi$ deux fonctions simples mesurables positives et
$\alpha, \beta \geq 0$. Alors~:
\begin{enumerate}
    \item (Lin\'earit\'e) $\int (\alpha\varphi + \beta\psi)\, d\mu =
    \alpha \int \varphi\, d\mu + \beta \int \psi\, d\mu$.
    \item (Monotonie) Si $\varphi \leq \psi$, alors
    $\int \varphi\, d\mu \leq \int \psi\, d\mu$.
    \item Si $\varphi = 0$ $\mu$-p.p., alors $\int \varphi\, d\mu = 0$.
    \item Pour tout $A \in \mathcal{A}$,
    $\int_A \varphi\, d\mu := \int_X \varphi \cdot \mathbf{1}_A\, d\mu$.
\end{enumerate}
\end{proposition}

\begin{proof}
(1) Consid\'erons la partition commune engendr\'ee par les niveaux de $\varphi$ et
$\psi$. Si $\varphi = \sum_i a_i \mathbf{1}_{B_i}$ et $\psi = \sum_j b_j
\mathbf{1}_{C_j}$ avec $\{B_i\}$ et $\{C_j\}$ des partitions, alors les
$B_i \cap C_j$ forment une partition plus fine et
\[
\alpha\varphi + \beta\psi = \sum_{i,j} (\alpha a_i + \beta b_j)
\mathbf{1}_{B_i \cap C_j}.
\]
L'\'egalit\'e suit par calcul direct.

(2) On \'ecrit $\psi - \varphi \geq 0$ et on utilise la lin\'earit\'e.
\end{proof}

\section{Int\'egrale des fonctions mesurables positives}

\begin{definition}[Int\'egrale d'une fonction positive]
Soit $f : X \to [0, +\infty]$ mesurable. On d\'efinit~:
\[
\int_X f\, d\mu = \sup\left\{\int_X \varphi\, d\mu :
\varphi \text{ simple mesurable}, 0 \leq \varphi \leq f\right\}.
\]
\end{definition}

\begin{proposition}[Propri\'et\'es fondamentales]
Soient $f, g : X \to [0, +\infty]$ mesurables.
\begin{enumerate}
    \item (Monotonie) Si $f \leq g$, alors $\int f\, d\mu \leq \int g\, d\mu$.
    \item (Homog\'en\'eit\'e) Pour $c \geq 0$, $\int cf\, d\mu = c \int f\, d\mu$.
    \item $\int f\, d\mu = 0$ si et seulement si $f = 0$ $\mu$-p.p.
    \item Si $\mu(A) = 0$, alors $\int_A f\, d\mu = 0$.
\end{enumerate}
\end{proposition}

\begin{proof}
(3) Si $f = 0$ p.p., toute fonction simple $0 \leq \varphi \leq f$ v\'erifie
$\varphi = 0$ p.p., donc $\int \varphi\, d\mu = 0$.

R\'eciproquement, si $\int f\, d\mu = 0$, posons $A_n = \{f > 1/n\}$. Alors
$\frac{1}{n} \mathbf{1}_{A_n} \leq f$, donc
$\frac{1}{n} \mu(A_n) \leq \int f\, d\mu = 0$, d'o\`u $\mu(A_n) = 0$ pour tout $n$.
Comme $\{f > 0\} = \bigcup_n A_n$, on a $\mu(\{f > 0\}) = 0$.
\end{proof}

\begin{theoreme}[Tchebyshev / In\'egalit\'e de Markov]
\label{thm:markov}
Soit $f : X \to [0, +\infty]$ mesurable. Pour tout $a > 0$~:
\[
\mu(\{f \geq a\}) \leq \frac{1}{a} \int_X f\, d\mu.
\]
\end{theoreme}

\begin{proof}
$a \cdot \mathbf{1}_{\{f \geq a\}} \leq f$, donc
$a \cdot \mu(\{f \geq a\}) = \int a \cdot \mathbf{1}_{\{f \geq a\}}\, d\mu
\leq \int f\, d\mu$.
\end{proof}

\section{Additivit\'e de l'int\'egrale (fonctions positives)}

\begin{theoreme}
\label{thm:additivite_positive}
Soient $f, g : X \to [0, +\infty]$ mesurables. Alors~:
\[
\int_X (f + g)\, d\mu = \int_X f\, d\mu + \int_X g\, d\mu.
\]
\end{theoreme}

Ce r\'esultat est une cons\'equence du th\'eor\`eme de convergence monotone
(chapitre~\ref{ch:convergence}). On peut aussi le d\'emontrer directement par
approximation avec des fonctions simples.

\section{Int\'egrale des fonctions int\'egrables}

\begin{definition}[Fonction int\'egrable]
Une fonction mesurable $f : X \to \overline{\R}$ est dite \emph{int\'egrable}
(ou \emph{$\mu$-int\'egrable}, ou \emph{sommable}) si~:
\[
\int_X \abs{f}\, d\mu < \infty.
\]
Dans ce cas, on d\'efinit~:
\[
\int_X f\, d\mu = \int_X f^+\, d\mu - \int_X f^-\, d\mu,
\]
o\`u $f^+ = \max(f, 0)$ et $f^- = \max(-f, 0)$. On note $L^1(X, \mathcal{A}, \mu)$
(ou $L^1(\mu)$) l'espace des fonctions int\'egrables.
\end{definition}

\begin{attention}[title=\textbf{La forme $\infty - \infty$ est interdite}]
La condition d'int\'egrabilit\'e $\int |f|\, d\mu < \infty$ garantit que les deux
int\'egrales $\int f^+\, d\mu$ et $\int f^-\, d\mu$ sont finies, \'evitant ainsi
la forme ind\'etermin\'ee $\infty - \infty$.
\end{attention}

\begin{theoreme}[Propri\'et\'es de l'int\'egrale]
\label{thm:proprietes_integrale}
Soient $f, g \in L^1(\mu)$ et $\alpha, \beta \in \R$.
\begin{enumerate}
    \item (Lin\'earit\'e) $\int (\alpha f + \beta g)\, d\mu =
    \alpha \int f\, d\mu + \beta \int g\, d\mu$.
    \item (In\'egalit\'e triangulaire)
    $\abs{\int f\, d\mu} \leq \int \abs{f}\, d\mu$.
    \item (Monotonie) Si $f \leq g$ p.p., alors
    $\int f\, d\mu \leq \int g\, d\mu$.
    \item Si $f = g$ p.p., alors $\int f\, d\mu = \int g\, d\mu$.
    \item Si $f \geq 0$ p.p. et $\int f\, d\mu = 0$, alors $f = 0$ p.p.
\end{enumerate}
\end{theoreme}

\begin{proof}
(2) On a $-\abs{f} \leq f \leq \abs{f}$, donc par monotonie
$-\int \abs{f}\, d\mu \leq \int f\, d\mu \leq \int \abs{f}\, d\mu$.
\end{proof}

\section{Comparaison avec l'int\'egrale de Riemann}

\begin{theoreme}
Soit $f : [a, b] \to \R$ une fonction born\'ee.
\begin{enumerate}
    \item Si $f$ est Riemann-int\'egrable, alors $f$ est Lebesgue-mesurable et
    les deux int\'egrales co\"incident~:
    \[
    (\mathcal{R})\!\int_a^b f(x)\, dx = \int_{[a,b]} f\, d\lambda.
    \]
    \item $f$ est Riemann-int\'egrable si et seulement si l'ensemble de ses points
    de discontinuit\'e est de mesure de Lebesgue nulle (th\'eor\`eme de Lebesgue
    sur l'int\'egrabilit\'e au sens de Riemann).
\end{enumerate}
\end{theoreme}

\begin{proof}[Id\'ee de la preuve]
(1) Les sommes de Darboux inf\'erieures et sup\'erieures sont des int\'egrales de
fonctions simples par rapport \`a $\lambda$. En passant \`a la limite sur les
partitions, on obtient l'\'egalit\'e.

(2) Soit $D$ l'ensemble des discontinuit\'es de $f$. On montre que les sommes de
Darboux sup\'erieures et inf\'erieures convergent vers la m\^eme valeur si et
seulement si $\lambda(D) = 0$.
\end{proof}

\begin{formules}[title=\textbf{Int\'egrale de Lebesgue --- R\'esum\'e}]
Pour $f : X \to [0, +\infty]$ mesurable~:
\[
\int_X f\, d\mu = \sup\left\{\sum_{k=1}^m c_k \mu(A_k) :
0 \leq \sum_k c_k \mathbf{1}_{A_k} \leq f\right\}.
\]
Pour $f$ g\'en\'erale int\'egrable~:
\[
\int_X f\, d\mu = \int_X f^+\, d\mu - \int_X f^-\, d\mu.
\]
\end{formules}

\section{La mesure $\nu_f$ associ\'ee \`a une fonction int\'egrable}

\begin{proposition}
Si $f \in L^1(\mu)$, $f \geq 0$, la fonction d'ensemble
$\nu_f(A) = \int_A f\, d\mu$ est une mesure finie sur $(X, \mathcal{A})$,
absolument continue par rapport \`a $\mu$.
\end{proposition}

\begin{proposition}[Continuit\'e absolue de l'int\'egrale]
Soit $f \in L^1(\mu)$. Pour tout $\varepsilon > 0$, il existe $\delta > 0$ tel que~:
\[
\mu(A) < \delta \implies \int_A \abs{f}\, d\mu < \varepsilon.
\]
\end{proposition}

\begin{proof}
Posons $f_n = \min(\abs{f}, n)$. Alors $f_n \uparrow \abs{f}$ et par convergence
monotone, $\int f_n\, d\mu \uparrow \int \abs{f}\, d\mu < \infty$. Fixons $n$ tel
que $\int (\abs{f} - f_n)\, d\mu < \varepsilon/2$. Alors pour $\mu(A) < \delta$
avec $\delta = \varepsilon/(2n)$~:
\[
\int_A \abs{f}\, d\mu \leq \int_A f_n\, d\mu + \int_A (\abs{f} - f_n)\, d\mu
\leq n \cdot \mu(A) + \varepsilon/2 < \varepsilon. \qedhere
\]
\end{proof}

\section{Exercices}

\begin{exercice}
Calculer $\int_{[0,1]} \mathbf{1}_{\Q}\, d\lambda$ et
$\int_{[0,1]} \mathbf{1}_{\R \setminus \Q}\, d\lambda$.
\end{exercice}

\begin{exercice}
Montrer que si $f \in L^1(\mu)$, alors $\{f \neq 0\}$ est $\sigma$-fini (pour $\mu$).
\end{exercice}

\begin{exercice}
Soit $f : [0, \infty) \to [0, \infty)$ mesurable. Montrer que
\[
\int_0^\infty f(x)\, dx = \int_0^\infty \lambda(\{f > t\})\, dt.
\]
(\emph{Formule de Cavalieri / formule de la couche.})
\end{exercice}

\begin{exercice}
Soit $f \in L^1(\R, \lambda)$. Montrer que
$\lim_{n \to \infty} \int_{\R} f(x) \sin(nx)\, dx = 0$
(\emph{lemme de Riemann--Lebesgue}).
\end{exercice}

\begin{exercice}
Soit $f \geq 0$ mesurable sur $(X, \mathcal{A}, \mu)$. Montrer que
$\nu(A) = \int_A f\, d\mu$ d\'efinit une mesure sur $(X, \mathcal{A})$ et que
pour toute $g \geq 0$ mesurable~:
$\int_X g\, d\nu = \int_X gf\, d\mu$.
\end{exercice}

\begin{exercice}
Soit $f : \R \to \R$ Riemann-int\'egrable sur tout $[a, b]$. Montrer que $f$ est
Lebesgue-mesurable. Est-elle n\'ecessairement dans $L^1(\R, \lambda)$~?
\end{exercice}
